
The results of this paper concern a single agent with no recursion and no formalized side
effects.
This section describes the limits of the prototype and what extending the theory requires.
Several of these extensions are subtle: designs that look natural turn out to be unsound, and we
note the pitfalls we know of.

\paragraph{The prototype.}
The oracle in our experiments is a seeded stand-in for an LLM, and the journal is in memory, so
the experiments show that the implementation follows the semantics, not how it behaves with
real models or persistent storage.
Experiments with real LLMs and a journal that survives process crashes are the natural next
step.

\paragraph{Recursion.}
Our terms have no recursion, so every run makes finitely many calls.
Agents that loop until a goal is met need recursion or a bounded loop; the theorems should
extend to finite prefixes of possibly infinite runs, but we have not proved this.

\paragraph{Side effects.}
Write-ahead logging of side effects is implemented and tested but not formalized.
A formal treatment must state what recovery can guarantee.
Performing an interrupted effect again is safe for an idempotent effect of a single agent, but
not in general: if another agent changed the same external state in the meantime, performing
the effect again can undo that change.
Compensation needs its own write-ahead discipline, or a crash during compensation can leave an
effect undone twice or not at all.
Li~\citep{li2026exactlyonce} shows that exactly-once effects cannot be achieved by verification
alone under late commits without a bound on in-flight time, so any recovery theorem must rest on
explicit assumptions about the external system, such as idempotence, atomicity, or a reliable
way to query whether an effect happened.

\paragraph{Concurrent agents.}
With several agents, replay and recovery become substantially harder.
\begin{itemize}[noitemsep]
  \item \emph{Interleaving.} A replayed run must follow the original interleaving of the agents,
    or agents that share state can read different values than they did originally. The
    scheduler's choices must therefore be journaled or reproduced deterministically.
  \item \emph{Recovery points.} Recovery must resume from a consistent cut of the agents'
    journals~\citep{chandy1985snapshots}. With shared memory this must account for which write
    each read observed, not only for messages; Wu et al.~\citep{wu2026safetoresume} show that
    agent frameworks can roll back to states that are not consistent cuts.
  \item \emph{Messages.} Messages sent before a cut but received after it must be delivered
    again exactly once, which requires identifying messages by channel and sequence number
    rather than by value.
  \item \emph{Restart.} Restarting part of a running system, as a supervisor might, also
    requires restoring the restarted agent's local state and suppressing messages it already
    sent.
\end{itemize}

\paragraph{Information-flow control.}
LLMbda~\citep{garby2026llmbda} proves probabilistic noninterference for a single agent's
program.
For concurrently scheduled agents, classical results on multi-threaded programs show that the
scheduler must not observe how long secret computations take, or timing becomes an information
channel~\citep{volpano1999probabilistic,sabelfeld2000probabilistic}.
Coarse-grained floating labels~\citep{stefan2011lio}, as used for agents by
APPA~\citep{kravchenko2026appa}, are a natural discipline, but they must keep the labels of
existing memory cells fixed: if a high-security agent can overwrite a public cell with a
secret label, whether a public agent can later read that cell leaks information.
Parallel composition inside one agent is similarly delicate, because one thread's secret
receive can block another's public output.
Peters-Gill et al.~\citep{petersgill2026camels} show empirically that single-agent defences
against prompt injection do not compose across agents, and Adams~\citep{adams2026sme} obtains
noninterference for language models by running one instance per security level; a
noninterference theorem for concurrently scheduled LLM agents remains open.

\paragraph{Compositional guarantees.}
If each agent of a system meets its own goal with probability at least $p_i$, Boole's
inequality bounds the probability that all do by $1 - \sum_i (1 - p_i)$, without independence.
The hard part is establishing the per-agent bounds inside a system, where an agent's inputs
depend on other agents and on the scheduler; this needs an assume-guarantee argument of the
kind Bensalem et al.~\citep{bensalem2026threelayer} advocate, and complements verification
approaches that bound the probability of violations directly~\citep{solkobreslin2026probverif}.
